% Compile with: uplatex file.tex && dvipdfmx file.dvi
%           or: latexmk -uplatex file.tex
%\documentclass[uplatex,dvipdfmx]{jsarticle}
\documentclass[dvipdfmx,a4paper,11pt]{jsarticle}
\usepackage{amsmath, amssymb}
\usepackage{tikz}

\title{scenario\_bcs\_hamiltonian}
\date{\today}

\begin{document}
\maketitle

BCSハミルトニアンを組み立てるための演算子・パラメータを準備する。$c$ は運動量 $k$・スピン $\sigma$（$u=\uparrow$, $d=\downarrow$）を持つ電子の消滅演算子。$g$ はこのあと Bogoliubov 変換で導入する準粒子の消滅演算子（特定の粒子種を持たないため、電子とは別に汎用フェルミオン演算子として宣言する）。$u, v$ は Bogoliubov 変換の係数（後で $u^2+v^2=1$ などの条件を課す）。$\xi_k$（xi）は電子の運動エネルギー（化学ポテンシャルからの相対値）、$\Delta$（Delta）はクーパー対を結びつけるペアポテンシャル（超伝導ギャップ）を表すスカラーパラメータ。

c を電子の消滅演算子として定義して

\[
  c
\]

g をフェルミオン演算子として定義して

\[
  g
\]

u, v をスカラーパラメータとして定義して

\[
  u, v
\]

xi, Delta をスカラーパラメータとして定義して

\[
  \xi, \Delta
\]

平均場BCSハミルトニアンを登録する。第1項 $\xi_k(c^\dagger_{k\uparrow}c_{k\uparrow}+c^\dagger_{-k\downarrow}c_{-k\downarrow})$ は電子の運動エネルギー、第2項 $\Delta(c^\dagger_{k\uparrow}c^\dagger_{-k\downarrow}+c_{-k\downarrow}c_{k\uparrow})$ はクーパー対の生成・消滅を表すペアポテンシャル項。

$H_BCS = \xi (c'(k,u) c(k,u) + c'(-k,d) c(-k,d)) + \Delta (c'(k,u) c'(-k,d) + c(-k,d) c(k,u))$

\begin{equation}
  H_BCS = \xi (c^{\dagger}(k,u) c(k,u) + c^{\dagger}(-k,d) c(-k,d)) + \Delta (c^{\dagger}(k,u) c^{\dagger}(-k,d) + c(-k,d) c(k,u))
  \label{eq:E1}
\end{equation}

運動方程式（ハイゼンベルク方程式 $dc/dt=i[H,c]$）を $c(k,u)$ について求め、この演算子が H\_BCS とどのように結びついているかを確認する。

c(k,u) の Heisenberg 方程式 E1

\begin{equation}
  \frac{d c(k,u)}{d t} = - \Delta \mathit{i} {c}_{-k,d}^{\dagger} - \xi \mathit{i} {c}_{ku}
  \label{eq:E2}
\end{equation}

Bogoliubov 変換（準粒子への基底変換）を、電子演算子とその随伴すべてを含む Nambu スピノル形式の4本まとめて定義する（ダガーは自動生成しない設計のため、生成演算子側の式も明示的に書く）。

$g(k,u) = u c(k,u) - v c'(-k,d)$

\begin{equation}
  g(k,u) = u c(k,u) - v c^{\dagger}(-k,d)
  \label{eq:E3}
\end{equation}

$g(-k,d) = u c(-k,d) + v c'(k,u)$

\begin{equation}
  g(-k,d) = u c(-k,d) + v c^{\dagger}(k,u)
  \label{eq:E4}
\end{equation}

$g'(k,u) = u c'(k,u) - v c(-k,d)$

\begin{equation}
  g^{\dagger}(k,u) = u c^{\dagger}(k,u) - v c(-k,d)
  \label{eq:E5}
\end{equation}

$g'(-k,d) = u c'(-k,d) + v c(k,u)$

\begin{equation}
  g^{\dagger}(-k,d) = u c^{\dagger}(-k,d) + v c(k,u)
  \label{eq:E6}
\end{equation}

先に定義した4本の変換式を、電子演算子 c(k,u), c(-k,d), c'(k,u), c'(-k,d) について解き直し（逆変換）、準粒子演算子gで電子演算子を表す式を得る。

E3, E4, E5, E6 を c(k,u), c(-k,d), c'(k,u), c'(-k,d) について逆変換して

\begin{equation}
  c_{k,u} = \frac{g^{\dagger}_{-k,d} \cdot v + g_{k,u} \cdot u}{u^{2} + v^{2}}
  \label{eq:E7}
\end{equation}

\begin{equation}
  c_{-k,d} = \frac{ - g^{\dagger}_{k,u} \cdot v + g_{-k,d} \cdot u}{u^{2} + v^{2}}
  \label{eq:E8}
\end{equation}

\begin{equation}
  c^{\dagger}_{k,u} = \frac{g^{\dagger}_{k,u} \cdot u + g_{-k,d} \cdot v}{u^{2} + v^{2}}
  \label{eq:E9}
\end{equation}

\begin{equation}
  c^{\dagger}_{-k,d} = \frac{g^{\dagger}_{-k,d} \cdot u - g_{k,u} \cdot v}{u^{2} + v^{2}}
  \label{eq:E10}
\end{equation}

得られた逆変換の結果を BCS ハミルトニアンへ代入し、準粒子演算子 $g$ で表されたハミルトニアンを得る。対生成項 $g^\dagger_{-k\downarrow}g^\dagger_{k\uparrow}$ の係数がゼロになるよう $u, v$ を選ぶ（Bogoliubov 条件）と、ハミルトニアンは準粒子数演算子 $g^\dagger g$ のみの対角形になる。

E1 の c(k,u) に E7 を、c(-k,d) に E8 を代入する

\begin{equation}
  H_BCS = \left(2 \frac{v^{2} \xi}{\left( u^{2} + v^{2} \right)^{2}} + \frac{2 \Delta u v}{\left( u^{2} + v^{2} \right)^{2}}\right)  + \left( - \frac{u v \xi}{\left( u^{2} + v^{2} \right)^{2}} + \frac{ - u v \xi}{\left( u^{2} + v^{2} \right)^{2}} - \Delta \left( \frac{ - v^{2}}{\left( u^{2} + v^{2} \right)^{2}} + \frac{u^{2}}{\left( u^{2} + v^{2} \right)^{2}} \right)\right) {g}_{-k,d}^{\dagger} {g}_{ku}^{\dagger} + \left(\frac{u^{2} \xi}{\left( u^{2} + v^{2} \right)^{2}} - \frac{v^{2} \xi}{\left( u^{2} + v^{2} \right)^{2}} - \frac{2 \Delta u v}{\left( u^{2} + v^{2} \right)^{2}}\right) {g}_{-k,d}^{\dagger} {g}_{-k,d} + \left(\frac{u^{2} \xi}{\left( u^{2} + v^{2} \right)^{2}} + \frac{ - 2 \Delta u v}{\left( u^{2} + v^{2} \right)^{2}} + \frac{ - v^{2} \xi}{\left( u^{2} + v^{2} \right)^{2}}\right) {g}_{ku}^{\dagger} {g}_{ku} + \left(2 \frac{u v \xi}{\left( u^{2} + v^{2} \right)^{2}} + \Delta \left( \frac{ - v^{2}}{\left( u^{2} + v^{2} \right)^{2}} + \frac{u^{2}}{\left( u^{2} + v^{2} \right)^{2}} \right)\right) {g}_{-k,d} {g}_{ku}
  \label{eq:E11}
\end{equation}

\end{document}